\chapter{Code Listings}

OmniLaTeX provides syntax-highlighted code listings via the
\texttt{minted} package, configured in \texttt{omnilatex-listings.sty}.

\section{minted Setup}

The default style is \texttt{manni} with a left-line frame, line numbers
on every line, and caching enabled for fast recompilation:

\begin{verbatim}
\usemintedstyle{manni}
\setminted{
    frame      = lines,
    framesep   = 2mm,
    linenos,
    numbersep  = 8pt,
    cache      = true,
}
\end{verbatim}

The cache directory is derived from the job name so that parallel builds
do not collide.

\section{Inline Code}

Inline code fragments use \texttt{\textbackslash mintinline}:

\begin{verbatim}
The function \mintinline{python}{print("hello")} writes to stdout.
\end{verbatim}

The language is specified as the first mandatory argument.

\section{Block Listings}

Full listings are placed inside the \texttt{minted} environment:

\begin{verbatim}
\begin{minted}{python}
def greet(name):
    print(f"Hello, {name}!")
\end{minted}
\end{verbatim}

Line numbers and the frame are applied automatically through the global
\texttt{\textbackslash setminted} defaults.

\section{Long Listings}

For listings that span multiple pages, wrap the environment in a
\texttt{longlisting} float:

\begin{verbatim}
\begin{longlisting}[H]
    \inputminted{lua}{scripts/init.lua}
    \caption{Initialisation script}
    \label{lst:init}
\end{longlisting}
\end{verbatim}

\texttt{longlisting} is a float that may break across pages while
preserving the listing counter and caption numbering.

\section{Line Highlighting}

Specific lines can be highlighted with the \texttt{highlightlines} option:

\begin{verbatim}
\begin{minted}[highlightlines={3,7-9}]{python}
x = 42          # highlighted
y = x + 1
z = y * 2       # highlighted
\end{minted}
\end{verbatim}

OmniLaTeX also defines placeholder commands for annotating template code:

\begin{table}[htbp]
    \centering
    \caption{Placeholder highlight commands}
    \begin{tabular}{@{} l l @{}}
        \toprule
        \textbf{Command} & \textbf{Colour} \\
        \midrule
        \verb|\phstring{text}| & red    \\
        \verb|\phnum{text}|    & blue   \\
        \verb|\phother{text}|  & green  \\
        \verb|\phnote{text}|   & orange \\
        \bottomrule
    \end{tabular}
\end{table}

These are intended for documentation that shows replaceable portions of
code.

\section{Language Support}

\texttt{minted} leverages Pygments, so any Pygments lexer is available.
Languages commonly used with OmniLaTeX include:

\begin{multicols}{2}
\begin{itemize}
    \item Python
    \item MATLAB
    \item Lua
    \item Modelica
    \item C / C++
    \item bash
    \item Dockerfile
    \item JSON
    \item SQL
    \item YAML
\end{itemize}
\end{multicols}

Add the \texttt{-{}-shell-escape} flag when invoking the compiler so that
Pygments can be called.

\section{Simulink Icons}

\texttt{omnilatex-graphics.sty} provides \verb|\mtlbsmlkicon{name}| for
embedding Simulink block icons inline:

\begin{verbatim}
\mtlbsmlkicon{Gain}
\end{verbatim}

The icon is rendered at a fixed size suitable for inline use.

\section{Accessible Line Numbers}

For PDF/UA-1 compliance, line-number digits are hidden from screen readers
via

\begin{verbatim}
\emptyaccsupp
\end{verbatim}

This macro is applied automatically inside every \texttt{minted}
environment when the accessibility mode is active.
